\section*{Local accuracy, geometry and the evaluation contract}
In the zonal loss

\[\mathcal L=\alpha\,\frac{1}{N_v}\sum_{i\in V}e_i^2
+(1-\alpha)\,\frac{1}{N_m}\sum_{i\in M}e_i^2,\]

$V$ denotes reverse-flow samples, $M$ the remaining samples, and $N_v,N_m$
their counts. The scalar $\alpha$ weights regional mean errors, rather than
the number of pixels in each region. Select it on validation cases, then freeze
it before opening final cases. Always show unweighted global AND reverse-flow
validation errors beside the weighted fits. A validation guardrail, such as a
two-percentage-point global penalty, is a selection rule, not a guarantee on new
geometries. Map the largest local error; it may sit at the step lip rather than
inside the reverse-flow zone.

The classroom coordinate model is an MLP. A strict branch/trunk model that uses
only $(x/H,y/H)$ in its trunk has a different geometry representation and training
problem. A worse result is not by itself proof of an architectural limitation.
Compare seeds, budget, training error and explicit wall-relative features before
attributing the difference to geometry. Do not mix numbers from different runs.

In nozzle registration, a station returned only at grid nodes cannot distinguish
sub-cell shock motion. Equal detected stations do not prove identical physical
shocks. Report grid spacing, ties and every non-monotone pair. Target-assisted
alignment is a representation audit; deployment must use a training-only locator.
Historical holdouts are regression cases once inspected, not new blind cases.

Branch sensors and query coordinates have different roles. Fixed branch sensors
describe the input consistently across cases. Each query must still be paired
with its own target value; equal array shapes alone do not establish that pairing.
